% source: David Mumford, "Abelian Varieties", TIFR Studies in Mathematics 5, Oxford UP, 1970
% pdf_page: 0066
% doc_page: 55 (Chapter II, end of Section 5 "Cohomology and base change";
%           beginning of Section 6, "The theorem of the cube: I")
% retrieved: 2026-07-11
% transcribed_by: claude-fable-5 (vision, rendered at 200dpi; scanned book, no text layer)

% requires: amsmath, amssymb

\DeclareMathOperator{\Ima}{Im}
\DeclareMathOperator{\Ker}{Ker}
\DeclareMathOperator{\Pic}{Pic}

% [running head: ALGEBRAIC THEORY VIA VARIETIES, page 55]

% [conclusion of the proof of Corollary 6 (Seesaw Theorem---provisional form)
%  from the previous page]
is an isomorphism. It clearly follows from the triviality of
$L \mid X \times \{t\}$ that the natural map
$p_2^*(\underline{M}) \to \underline{L}$ is an isomorphism. Since
$\underline{M}$ is the sheaf of sections of $M$, then $p_2^* M \simeq L$.

\section*{6. The theorem of the cube: I}

\paragraph{Theorem.}
\textit{Let $X$, $Y$ be complete varieties, $Z$ any variety and $x_0$, $y_0$
and $z_0$ base points on $X$, $Y$, and $Z$, respectively. If $L$ is any line
bundle on $X \times Y \times Z$ whose restrictions to each of
$\{x_0\} \times Y \times Z$, $X \times \{y_0\} \times Z$ and
$X \times Y \times \{z_0\}$ are trivial, $L$ is trivial.}

\paragraph{Remark.}
Let $T$ be a contravariant functor on the category of
% [the print shows a stray prime after this first "T"; the functor is called T throughout]
complete varieties into the category $\underline{\mathrm{Ab}}$ of abelian
groups. Let $X_0, \ldots, X_n$ be any system of complete varieties, $x_i^0$ a
base point of $X_i$, and let
$\pi_i : X_0 \times \ldots \times X_n \to
 X_0 \times \ldots \times \hat{X}_i \times \ldots \times X_n$
($\hat{X}_i$ indicating the omission of the $i$-th factor $X_i$) be the
projection map, and
$\sigma_i : X_0 \times \ldots \times \hat{X}_i \times \ldots \times X_n \to
 X_0 \times \ldots \times X_n$
the `inclusion' defined by
\[
  \sigma_i(x_0, \ldots, x_{i-1}, x_{i+1}, \ldots, x_n) =
  (x_0, \ldots, x_{i-1}, x_i^0, x_{i+1}, \ldots, x_n).
\]
Consider the homomorphisms
\[
  \alpha_T^n : \prod_{i=0}^{n}
  T(X_0 \times \ldots \times \hat{X}_i \times \ldots \times X_n) \to
  T(X_0 \times \ldots \times X_n),
\]
\[
  \beta_T^n : T(X_0 \times \ldots \times X_n) \to \prod_{i=1}^{n}
  T(X_0 \times \ldots \times \hat{X}_i \times \ldots \times X_n)
\]
defined by
\[
  \alpha_T^n(\xi_0, \ldots, \xi_n) = \sum_0^n \pi_i^*(\xi_i), \quad
  \beta_T^n(\eta) = (\sigma_0^*(\eta), \sigma_2^*(\eta), \ldots, \sigma_n^*(\eta)).
\]
% [sic: the components of beta are printed with subscripts 0, 2, ..., n,
%  while the product defining its target runs over i = 1, ..., n]
One then proves by an easy induction on $n$ that we have a natural splitting
$T(X_0 \times \ldots \times X_n) = \Ima \alpha \oplus \Ker \beta$. The
functor $T$ is said to be of \textit{order} $n$ (\textit{linear} if $n = 1$,
\textit{quadratic} if $n = 2$, etc.) if $\alpha$ is surjective, or
equivalently $\beta$ is injective. (Note that the definition of $\alpha$ is
independent of base points.)

Thus, the above theorem (when $Z$ is also assumed complete) may be
paraphrased as saying that the functor $\Pic X$ is a quadratic functor on the
category of complete varieties.
